\section{Introduction}
\label{sec:intro}

% 11-07：按 07 大纲与 11-05 方法写。三条贡献：界与时间界、由界导出的方法、实验。
% 11-07 审核改：去掉修饰词与无法兑现的说法（"certified" 改为按长期口径检查）；缩写首次出现时定义；引言不提前用"相位"记号。

Inspection, patrol and security robots keep many sites staffed around the clock.
Each robot serves a duty at a site, returns to its depot and is maintained before it leaves again.
Maintenance follows a cycle: a few short services, then an overhaul that returns the robot to its initial state.
An operator buys the fleet once and rotates it so that every site is staffed in every slot for years.
The question is how many robots and tools to buy and how to rotate them.

We call this problem multi-domain rotation rostering (MDRR): several sites served from one depot, several robot types with different maintenance cycles, shared tools, and an infinite horizon.
Three features make it hard.
First, the horizon is long.
The cycles of different robot types do not divide each other, so a good schedule may repeat only after years, and one period is too long to encode and search.
Second, a new fleet starts in step.
All robots start fresh on the same day, their first overhauls fall close together, and extra robots are needed to cover them.
Steady-state models do not see this start-up cost, and it can grow linearly with the number of posts.
Third, a schedule that covers a finite horizon can run short later, so a finite-horizon evaluation can accept infeasible schedules.

Existing work covers parts of this problem.
Cyclic staffing, crew rostering and rolling stock rotation planning optimize a steady state in which the fleet is already spread over its cycles~\cite{bartholdi1980cyclic,breugem2023threephase,wang2024rollingstock,feng2024branchprice}.
Persistent robot service has been studied with recharging~\cite{asghar2025multirobot,arribas2023uav,kumar2025persistent}, and synchronized battery depletion breaks steady-state fleet sizing~\cite{ramesh2026robust}.
Fluid relaxations give lower bounds and asymptotically optimal policies for large scheduling and resource allocation systems~\cite{bertsimas2002fluid,brown2025fluid,gast2024lp}, and a fluid allocation can be turned into an exactly feasible repeating schedule~\cite{zheng2026highmultiplicity}.
Evolutionary algorithms with direct or random-key encodings need genotypes that grow with the number of decisions~\cite{londe2025brkga,ding2024hfmae}.
Dispatching rules evolved by genetic programming~\cite{zhang2024survey,xu2024lexicase}, trained by reinforcement learning~\cite{song2023fjsp,liu2025trainer} or generated by large language models~\cite{liu2024eoh,ye2024reevo,yu2026lsh} scale to long horizons, but they give no lower bound and are evaluated over a finite horizon.

This paper brackets the optimum of MDRR and builds the method from the bracket.
On the lower side, a rate relaxation assumes that maintenance is spread evenly over the robots, and a start-up relaxation adds the fresh start.
On the upper side, an entering theorem lets a new fleet join any feasible rotation: each new robot waits until the robot it copies finishes its first overhaul, and a few extra robots serve the duties before that.
After the start-up, the schedule repeats the rotation, so the search can focus on the start-up; Section~\ref{sec:experiments} shows that a window of about one week is enough.
The method follows the bracket step by step.
The rate relaxation gives a rotation rule, the entering theorem gives the start-up, and an evolutionary algorithm changes a few decisions in a short start-up window and lowers the fleet.
Every returned schedule passes a long-run feasibility check.

The contributions are as follows.
\begin{enumerate}
  \item Bounds and a time bound for MDRR.
  The rate and start-up relaxations give lower bounds.
  For one robot type, an even rotation attains the rate bound, the start-up cost without early overhaul grows linearly with the number of posts, and early overhaul removes it when the turnaround times meet a simple condition.
  The entering theorem gives an upper bound and a start-up window whose length does not depend on the numbers of sites, posts and robot types.
  \item A method derived from these results, rule-deviation decoding on a constructed start-up (C-RDD).
  It reads a rotation rule from the rate relaxation, builds the start-up by the entering theorem and encodes only a few deviations inside a short start-up window, so the genotype length does not depend on the horizon.
  \item Experiments that compare C-RDD with dispatching rules, other decoders and a rolling-horizon integer program under the same infinite-horizon evaluation, and ablations of the rule, the construction, the search and the evaluation.
\end{enumerate}

Section~\ref{sec:related} reviews related work.
Section~\ref{sec:problem} defines MDRR.
Section~\ref{sec:bounds} derives the bounds and the time bound.
Section~\ref{sec:method} presents C-RDD, Section~\ref{sec:experiments} reports the experiments, Section~\ref{sec:discussion} discusses the remaining gap, and Section~\ref{sec:conclusion} concludes.
